% ============================================================================
% Fichier  : partie4/P04_C13_grc_pratique.tex
% Titre    : La GRC en Pratique
% Auteur   : MINKA MI NGUIDJOÏ Thierry Emmanuel
% Version  : 1.0 -- novembre 2026
% Corrigés : corriges/Corr_C13.tex
% ============================================================================
% Faits du cas (bible) : BANQUE SANAGA, registre (R-AG-01, R-AG-02, R-DC-02,
% R-DC-03, R-DC-06), plan de traitement, audits de 2028, CERTIQUAL,
% indicateurs I1 à I6 (I6 = 11 %), DIGICORE (rapport ISAE 3402 de type 2),
% avenant GAB, monétique sous PCI DSS, règle « temporaire » de Garoua chez
% LiZbiZ (leçon reprise).
% Faits nouveaux (reportés dans la bible) : M. ATANGANA, responsable GRC
% nommé en octobre 2028, rattaché au directeur des risques ; comité des
% risques de janvier 2029 ; registre des obligations O1 à O12 (7/2/3) ;
% matrice risques-contrôles C1 à C9 ; indicateurs de risque K1 à K5 ;
% contrôles automatisés quotidiens ; DIGICORE n'a pas remis son rapport à
% l'échéance de décembre 2028.
% Loi n° 2010/012 : art. 26, 29, 30 résumés (Chapitre 10).
% ============================================================================

\chapter{La GRC en Pratique}
\label{chap:grc-pratique}
\chaptermark{La GRC en pratique}

\epigraph{%
    « Quatre rapports qui ne se parlent pas\\
    valent moins qu'un tableau qui les relie. »%
}{Principe du cours}

\niveauChapitre{Intermédiaire}{%
    Maîtrise du Chapitre~\ref{chap:smsi}%
}

\begin{boxobjectifs}
\begin{itemize}
    \item Relier risques, obligations, contrôles et preuves dans un
          dispositif unique.
    \item Tenir un registre des risques et un registre des obligations.
    \item Construire une matrice risques-contrôles, et y repérer les risques
          sans contrôle et les contrôles sans risque.
    \item Distinguer indicateurs de risque et indicateurs de performance, et
          fixer des seuils.
    \item Concevoir le tableau de bord GRC d'une instance de gouvernance, y
          compris ce que l'on ignore.
    \item Organiser le contrôle continu de la conformité, et en connaître
          les limites.
\end{itemize}
\end{boxobjectifs}

\bigskip

% ============================================================================
\section*{Avant de commencer : quatre rapports}
% ============================================================================

Douala, janvier 2029. Le comité des risques du conseil de la BANQUE
SANAGA\footnote{Comme dans tout le cours, la BANQUE SANAGA, ses dirigeants,
ses prestataires, les dates et les montants sont fictifs~; les textes cités
sont exacts.} reçoit, pour la même séance, quatre documents~: le tableau de
bord de la RSSI, le rapport annuel de l'audit interne, le rapport de
CERTIQUAL AFRIQUE et une lettre de la COBAC. Chacun est juste, et aucun ne
parle la langue des autres. Le tableau de bord dit «~100~\% des comptes avec
second facteur~»~; l'audit interne, «~trois non-conformités sur les
agences~»~; CERTIQUAL, «~une majeure levée~»~; la COBAC demande «~l'état des
mesures prises depuis l'incident de 2026~».

Le président du comité pose une seule question~: «~Où en sommes-nous~?~»
M.~ATANGANA, responsable GRC nommé en octobre~2028 et rattaché au directeur
des risques, a trois mois pour qu'une page y réponde.

\medskip

Ce chapitre montre comment faire vivre la gouvernance, le risque et la
conformité dans un seul dispositif~: des registres reliés, une matrice, des
indicateurs, un tableau de bord, et des contrôles qui se vérifient seuls.

\begin{boxalert}
\textbf{Violation typique.} Présenter à la gouvernance des rapports
juxtaposés (sécurité, audit, conformité, certification), chacun avec ses
références et ses échelles.

\smallskip
\textbf{Règle~:} un même référentiel relie risques, obligations, contrôles
et preuves. Chaque rapport est une vue de ce référentiel, et le tableau de
bord de la gouvernance les réunit.
\end{boxalert}

\separateur

% ============================================================================
\section{Un Dispositif Unique}
\label{sec:c13-dispositif}
% ============================================================================

Le Chapitre~\ref{chap:grc} a défini la GRC comme un seul dispositif, avec
un référentiel commun. En pratique, ce référentiel tient en quatre objets
et leurs liens.

\begin{figure}[H]
\centering
\begin{tikzpicture}[
    objet/.style={draw=CouleurPrimaire, thick, rounded corners,
                  fill=CouleurPrimaire!8, text width=2.6cm,
                  minimum height=1.1cm, align=center, font=\small},
    lien/.style={<->, thick, CouleurSecondaire},
    etiq/.style={font=\scriptsize, fill=white, inner sep=1pt}]
    \node[objet] (R) at (0,0) {\textbf{Risques}\\registre};
    \node[objet] (O) at (0,-3) {\textbf{Obligations}\\registre};
    \node[objet] (C) at (4.5,-1.5) {\textbf{Contrôles}\\base commune};
    \node[objet] (P) at (9,-1.5) {\textbf{Preuves}\\tests, extractions};
    \draw[lien] (R) -- node[etiq]{traité par} (C);
    \draw[lien] (O) -- node[etiq]{satisfaite par} (C);
    \draw[lien] (C) -- node[etiq]{démontré par} (P);
\end{tikzpicture}
\caption{Le référentiel commun de la GRC}
\label{fig:c13-referentiel}
\end{figure}

Chaque objet a son propriétaire~: le risque, un directeur métier~;
l'obligation, la fonction conformité ou juridique~; le contrôle, son
exploitant~; la preuve, celui qui l'a produite et celui qui l'a évaluée.
Les liens sont le travail propre du responsable GRC~: c'est par eux qu'un
constat d'audit change la cote d'un risque, ou qu'une nouvelle loi crée un
contrôle.

% ============================================================================
\section{Les Deux Registres}
\label{sec:c13-registres}
% ============================================================================

\subsection{Le registre des risques}

Le registre des risques reprend, pour chaque risque, les rubriques des
chapitres précédents~: scénario, propriétaire, cotes inhérente et
résiduelle avec leur base de preuve (Chapitre~\ref{chap:risque}), décision
de traitement ou d'acceptation (Chapitre~\ref{chap:iso27005}), date de
revue.

\begin{tableaucours}{L{1.5cm}L{3.4cm}L{2.1cm}cL{3.0cm}}{%
    \tableauHeader{Réf.} &
    \tableauHeader{Risque} &
    \tableauHeader{Propriétaire} &
    \tableauHeader{Rés.} &
    \tableauHeader{Décision, revue}
}
R-AG-01 & Rançongiciel en agence & Dir. réseau d'agences & 6 & Traité~;
revue juin 2029 \\
R-AG-02 & Compromission d'un GAB & Dir. réseau d'agences & 9 &
Remplacement en cours~; revue mars 2029 \\
R-DC-02 & Défaillance de DIGICORE & Dir. des opérations & 6 (\NEK{}) &
Accepté en l'état des connaissances~; rapport attendu \\
R-DC-03 & Intrusion par Internet & DSI & 8 & Traité~; revue juin 2029 \\
R-DC-06 & Déni de service des services en ligne & Dir. canaux numériques &
9 & Protection contractualisée~; test en mars \\
\end{tableaucours}
\vspace{-0.8em}
\begin{center}{\small\itshape Tableau~13.1. Extrait du registre des risques
(janvier 2029)}\end{center}

\subsection{Le registre des obligations}

\begin{boxdef}
\termecle{Registre des obligations}\index{registre des obligations}~: liste
des exigences qui s'imposent à l'organisation (lois, règlements, contrats,
normes auxquelles elle s'engage), avec, pour chacune, sa source, son
contenu résumé, les contrôles qui la satisfont, la preuve et son statut.
\end{boxdef}

\begin{tableaucours}{cL{3.6cm}L{4.4cm}c}{%
    \tableauHeader{Réf.} &
    \tableauHeader{Source} &
    \tableauHeader{Exigence (résumée)} &
    \tableauHeader{Statut}
}
O1 & Loi n°~2010/012, art.~26~(1) & Gestion continue des risques & \OUI{} \\
O2 & Loi n°~2010/012, art.~26~(2)--(3) & Mécanismes de sécurité visés par
l'Agence & \NEK{} \\
O3 & Loi n°~2010/012, art.~29~(1) & Conservation dix ans des données de
connexion & \NON{} \\
O4 & Loi n°~2010/012, art.~29~(2) & Surveillance des accès aux données &
\OUI{} \\
O5 & Loi n°~2010/012, art.~30~(1) & Revue périodique des systèmes de
sécurité & \OUI{} \\
O6 & Règlement COBAC R-2016/04 & Contrôle interne couvrant le risque
opérationnel & \OUI{} \\
O7 & PCI DSS (monétique) & Évaluation annuelle de conformité & \NEK{} \\
O8 & ISO/IEC~27001 (certificat) & Maintien du SMSI certifié & \OUI{} \\
O9 & Contrat DIGICORE & Rapport d'assurance annuel de type~2 & \NON{} \\
O10 & Avenant du prestataire des GAB & Droit d'audit de la banque & \OUI{} \\
O11 & Loi n°~2024/017 (à vérifier) & Obligations sur les données
personnelles & \NEK{} \\
O12 & Contrats des correspondants bancaires & Questionnaire de sécurité
annuel & \OUI{} \\
\end{tableaucours}
\vspace{-0.8em}
\begin{center}{\small\itshape Tableau~13.2. Le registre des obligations
(janvier 2029)}\end{center}

\begin{boxlecture}{L'état des obligations}
\textbf{Constatation.} 7~\OUI{}, 2~\NON{}, 3~\NEK{} sur 12~:
\[
\mu \approx 0{,}58~;\quad U = 0{,}25~;\quad
H = -\left(\tfrac{7}{12}\log_2\tfrac{7}{12} + \tfrac{2}{12}\log_2\tfrac{2}{12}
  + \tfrac{3}{12}\log_2\tfrac{3}{12}\right) \approx 0{,}454 + 0{,}431 + 0{,}500
  \approx 1{,}38~\text{bit}.
\]

\textbf{Analyse} (confiance élevée). Conformité opaque~: la banque satisfait
la majorité de ses obligations, mais un quart d'entre elles est dans
l'inconnu. Les trois \NEK{} n'ont pas la même cause~: O2 est une pièce à
retrouver~; O7 attend le rapport de l'évaluateur de la monétique~; O11
dépend d'un texte que la banque n'a pas encore analysé. Les deux \NON{}
sont connus et datés~: O3 tient à la capacité de stockage (deux ans au lieu
de dix), O9 au retard de DIGICORE, qui n'a pas remis son rapport à
l'échéance de décembre 2028.

\medskip
\hfill\lecture{0{,}58}{0{,}25}{1{,}38}
\end{boxlecture}

% ============================================================================
\section{La Matrice Risques-Contrôles}
\label{sec:c13-matrice}
% ============================================================================

\subsection{Croiser}

\begin{boxdef}
\termecle{Matrice risques-contrôles}\index{matrice risques-contrôles}~:
tableau qui croise les contrôles de la base commune et les risques du
registre, en indiquant quel contrôle traite quel risque, et avec la dernière
cote de chaque contrôle.
\end{boxdef}

\begin{tableaucours}{cL{3.6cm}cccccc}{%
    \tableauHeader{Réf.} &
    \tableauHeader{Contrôle} &
    \tableauHeader{AG-01} & \tableauHeader{AG-02} &
    \tableauHeader{DC-02} & \tableauHeader{DC-03} &
    \tableauHeader{DC-06} & \tableauHeader{Cote}
}
C1 & Second facteur & $\times$ & & & $\times$ & & \OUI{} \\
C2 & Cloisonnement & $\times$ & $\times$ & & & & \OUI{} \\
C3 & Sauvegardes testées & $\times$ & & $\times$ & & & \OUI{} \\
C4 & Mode dégradé & $\times$ & $\times$ & $\times$ & & $\times$ & \NON{} \\
C5 & Journalisation & $\times$ & $\times$ & & $\times$ & & \OUI{} \\
C6 & Correctifs & $\times$ & $\times$ & & $\times$ & & \NEK{} \\
C7 & Sensibilisation & $\times$ & & & $\times$ & & \NON{} \\
C8 & Assurance des prestataires & & $\times$ & $\times$ & & & \NON{} \\
C9 & Revue mensuelle des règles de pare-feu & & & & & & \OUI{} \\
\end{tableaucours}
\vspace{-0.8em}
\begin{center}{\small\itshape Tableau~13.3. Matrice risques-contrôles
(extrait)}\end{center}

\subsection{Lire la matrice}

La matrice se lit dans trois directions.
\begin{itemize}
    \item \textbf{Par ligne}~: les contrôles qui traitent beaucoup de risques
          sont des \emph{contrôles clés}~; C4 (mode dégradé) en traite
          quatre, et il est \NON{}~: c'est la priorité du plan.
    \item \textbf{Par colonne}~: un risque dont tous les contrôles sont
          \NON{} ou \NEK{} n'est pas maîtrisé. R-DC-06 n'a qu'un contrôle,
          C4, coté \NON{}~: la protection contre le déni de service,
          contractualisée avec le fournisseur d'accès, n'est pas encore un
          contrôle de la base~; il faut l'y inscrire et la tester.
    \item \textbf{Les lignes vides}~: C9 ne traite aucun risque du registre.
          Soit il est inutile, soit un risque manque. C'est le second cas~:
          la règle «~temporaire~» de LiZbiZ (Chapitre~\ref{chap:audit-technique})
          a montré le risque d'une règle de filtrage qui survit à son besoin~;
          le registre de la banque doit le porter.
\end{itemize}

\begin{boxprocedure}
\textbf{Tenir la matrice.} Le responsable GRC vérifie, à chaque mise à
jour~:
\begin{itemize}
    \item que chaque risque au-delà de l'appétence a au moins un contrôle
          coté \OUI{} par étape de son scénario~;
    \item que chaque contrôle est relié à au moins un risque ou une
          obligation~;
    \item que la cote de chaque contrôle vient du dernier test, avec sa
          date et sa source (audit, contrôle automatisé, autoévaluation)~;
    \item qu'un constat d'audit modifie la cote dans la semaine, et la cote
          du risque à la revue suivante.
\end{itemize}
\end{boxprocedure}

% ============================================================================
\section{Les Indicateurs de Risque}
\label{sec:c13-kri}
% ============================================================================

\subsection{Risque ou performance}

Les indicateurs du SMSI (Chapitre~\ref{chap:smsi}) mesurent la mise en
œuvre et l'efficacité des mesures. Un \termecle{indicateur de
risque}\index{indicateur de risque} mesure autre chose~: un signal qui
annonce que l'exposition change, souvent avant l'incident. Il se lit avec
des seuils, qui déclenchent une action.

\begin{tableaucours}{cL{4.2cm}cccL{2.6cm}}{%
    \tableauHeader{Réf.} &
    \tableauHeader{Indicateur de risque} &
    \tableauHeader{Vert} &
    \tableauHeader{Orange} &
    \tableauHeader{Rouge} &
    \tableauHeader{Risque suivi}
}
K1 & Agents qui cliquent lors d'une campagne simulée & $<$~5~\% & 5 à
10~\% & $>$~10~\% & R-AG-01 \\
K2 & Vulnérabilités exploitées non corrigées après 7~jours & 0 & 1 à 2 &
$\geq$~3 & R-AG-01, R-DC-03 \\
K3 & Règles de filtrage temporaires échues & 0 & 1 & $\geq$~2 & R-AG-01,
R-DC-03 \\
K4 & Prestataires critiques sans rapport d'assurance à jour & 0 & 1 &
$\geq$~2 & R-DC-02, R-AG-02 \\
K5 & Tentatives de connexion refusées sur l'accès distant (par semaine,
écart à la moyenne) & $<$~2 fois & 2 à 5 fois & $>$~5 fois & R-DC-03 \\
\end{tableaucours}
\vspace{-0.8em}
\begin{center}{\small\itshape Tableau~13.4. Les indicateurs de risque et
leurs seuils}\end{center}

\begin{boxinfo}
\textbf{Ce que déclenche un seuil.} Un indicateur au rouge ne change pas la
cote du risque à lui seul~: il déclenche la revue du risque par son
propriétaire, sans attendre la date prévue (Chapitre~\ref{chap:iso27005}).
En janvier 2029, K1 est au rouge (11~\%) et K4 à l'orange (DIGICORE)~: les
revues de R-AG-01 et de R-DC-02 sont avancées.
\end{boxinfo}

% ============================================================================
\section{Le Tableau de Bord GRC}
\label{sec:c13-tableau}
% ============================================================================

\subsection{Une page pour la gouvernance}

Le tableau de bord du comité des risques n'est pas celui de la RSSI. Il ne
détaille pas les mesures~; il répond à trois questions~: où sont les risques
qui dépassent l'appétence~? Respectons-nous nos obligations~? Que
savons-nous, et que ne savons-nous pas~?

\begin{boxscenario}{Opération SANAGA. La page de M. ATANGANA, avril 2029}
\textbf{1. Risques au-delà de l'appétence~:} R-AG-02 (GAB, 9~; remplacement
à 60~\%, fin prévue en juin)~; R-DC-06 (déni de service, 9~; test de la
protection en mars, réussi~: proposition de passer à 6).

\medskip
\textbf{2. Obligations~:} 7 satisfaites sur 12~; deux non satisfaites (O3,
conservation, budget de stockage demandé~; O9, DIGICORE, relance et
pénalité contractuelle)~; trois inconnues (O2, O7, O11), avec une date pour
chacune.

\medskip
\textbf{3. Contrôles clés~:} C4 (mode dégradé) non conforme, plan à fin
juin~; C6 (correctifs) inconnu, test de l'audit interne en mai.

\medskip
\textbf{4. Indicateurs de risque~:} K1 au rouge (11~\%)~; K4 à l'orange.

\medskip
\textbf{5. Ce que nous ne savons pas~:} $U = 0{,}25$ sur les obligations~;
contrôles C6 et sensibilisation à vérifier~; base \NEK{} de la vraisemblance
de R-DC-02.

\medskip
\textbf{6. Décisions demandées au comité~:} budget de stockage (O3)~;
position sur DIGICORE si le rapport n'est pas remis en juin (changement de
prestataire)~; validation de la nouvelle cote de R-DC-06.
\end{boxscenario}

La rubrique~5 est la plus rare dans les tableaux de bord, et la plus
utile. Un comité qui ne voit que des chiffres rassurants ne sait pas sur
quoi ils reposent~; un comité qui voit l'incertitude sait où porter
l'effort et ce qu'il accepte d'ignorer.

% ============================================================================
\section{Le Contrôle Continu de la Conformité}
\label{sec:c13-continu}
% ============================================================================

\subsection{Des contrôles qui se vérifient seuls}

Certains contrôles peuvent être testés automatiquement, chaque jour, par
une extraction comparée à une règle~: la part des comptes avec second
facteur (C1), la part des agences qui envoient des journaux (C5), les
règles de filtrage échues (C9, K3). C'est le \termecle{contrôle continu de
la conformité}\index{contrôle continu de la conformité}~: une alerte part
dès qu'une extraction s'écarte de la règle, au lieu d'attendre l'audit
suivant.

\begin{boxscenario}{Opération SANAGA. Trois tests quotidiens}
Depuis février 2029, trois scripts tournent chaque nuit~: extraction de
l'annuaire (comptes sans second facteur), extraction du collecteur
(agences silencieuses depuis plus de 24~heures), extraction des pare-feu
(règles dont la date d'expiration est passée). Leurs résultats alimentent
directement la cote de C1, C5 et C9 dans la matrice.

\medskip
\textbf{Le 14~mars}, le test de C5 signale une agence silencieuse depuis
trente heures~: un commutateur remplacé n'avait pas reçu la configuration
de référence. L'écart est corrigé le jour même. En 2028, le même défaut
avait attendu l'audit pour être découvert (Chapitre~\ref{chap:audit}).
\end{boxscenario}

\subsection{Les limites}

\begin{boxalert}
\textbf{Violation typique.} Tenir pour \OUI{} un contrôle parce que le test
automatique ne signale rien.

\smallskip
\textbf{Règle~:} un test automatique est lui-même un contrôle, qui doit
être validé~: voit-il toute la population (exhaustivité)~? Sa règle
est-elle juste (exactitude)~? Fonctionne-t-il encore (un script arrêté ne
signale rien)~? L'audit interne teste les tests, au moins une fois par an.
Un script qui ne tourne plus transforme un \OUI{} en \NEK{} sans que
personne ne le voie.
\end{boxalert}

Le contrôle continu de la conformité porte sur les contrôles~; l'audit
continu, qui porte sur les données et les transactions pour y détecter les
anomalies, est présenté au Chapitre~\ref{chap:emergents}.

\begin{boxinfo}
\textbf{La loi et le continu.} L'article~26 de la loi n°~2010/012 demande
aux exploitants de gérer \emph{continûment} les risques de sécurité
(Chapitre~\ref{chap:iso27005})~\cite{loi2010012}. Le dispositif décrit dans
ce chapitre (registres reliés, indicateurs à seuils, tests quotidiens) est
une manière concrète de donner un contenu à cet adverbe.
\end{boxinfo}

% ============================================================================
\section{Ce que la GRC demande à chaque Métier}
\label{sec:c13-metiers}
% ============================================================================

\begin{boxmetier}{Responsable GRC}
\textbf{Tenir les liens}
\begin{enumerate}
    \item Tenir le registre des obligations, avec la fonction juridique, et
          le relier aux contrôles.
    \item Tenir la matrice risques-contrôles à jour de chaque test et de
          chaque constat.
    \item Repérer les risques sans contrôle et les contrôles sans risque.
    \item Fixer avec les propriétaires les indicateurs de risque et leurs
          seuils.
    \item Produire une page pour la gouvernance, avec une rubrique «~ce que
          nous ne savons pas~».
\end{enumerate}
\end{boxmetier}

\begin{boxmetier}{RSSI}
\textbf{Alimenter le référentiel}

La RSSI ne tient plus son propre tableau à part~: ses indicateurs et ses
tests alimentent la matrice. Elle automatise les tests qui peuvent l'être,
et demande à l'audit interne de valider ses scripts.
\end{boxmetier}

\begin{boxmetier}{Auditeur en cybersécurité}
\textbf{Auditer le dispositif, et s'en servir}

L'auditeur peut s'appuyer sur la matrice et sur les tests automatiques pour
cibler sa mission, à condition de les avoir validés. Il audite aussi le
dispositif lui-même~: les liens sont-ils justes, les cotes à jour, les
seuils suivis~?
\end{boxmetier}

\begin{boxmetier}{Gestionnaire des risques cyber}
\textbf{Réagir aux seuils}

Un indicateur au rouge déclenche la revue du risque par son propriétaire.
Le gestionnaire des risques organise cette revue dans le mois, et en
inscrit la décision au registre.
\end{boxmetier}

\separateur

\begin{boxresume}
\textbf{Ce qu'il faut retenir de ce chapitre~:}
\begin{itemize}
    \item \textbf{Un référentiel commun}~: risques, obligations, contrôles,
          preuves, et leurs liens.
    \item \textbf{Deux registres}~: les risques (avec propriétaires et
          décisions) et les obligations (avec source et statut).
    \item \textbf{La matrice risques-contrôles} révèle les contrôles clés,
          les risques sans contrôle et les contrôles sans risque.
    \item \textbf{Les indicateurs de risque} annoncent un changement
          d'exposition~; leurs seuils déclenchent une revue.
    \item \textbf{Le tableau de bord de la gouvernance} tient en une page,
          et dit ce que l'on ne sait pas.
    \item \textbf{Le contrôle continu} détecte les écarts au jour le jour~;
          les tests automatiques sont eux-mêmes des contrôles à valider.
\end{itemize}
\end{boxresume}

\begin{boxinfo}
\textbf{Points de contrôle.} Vous avez compris ce chapitre si vous savez~:
\begin{itemize}
    \item rédiger une ligne du registre des obligations à partir d'un
          article de loi~;
    \item lire une matrice risques-contrôles dans ses trois directions~;
    \item définir un indicateur de risque et ses seuils~;
    \item dire pourquoi un test automatique silencieux n'est pas une
          preuve.
\end{itemize}
\end{boxinfo}

% ============================================================================
\section*{Exercices}
% ============================================================================

\begin{boxexo}[title={\ding{45}\quad Exercice~13.1. Classer \niveaubase}]
Dites si chacun des éléments suivants est un risque, une obligation, un
contrôle, une preuve, un indicateur de performance ou un indicateur de
risque.
\begin{enumerate}[(a)]
    \item Conserver dix ans les données de connexion.
    \item La part des comptes avec second facteur.
    \item Le procès-verbal de restauration du 12~février.
    \item Une défaillance prolongée de DIGICORE.
    \item Le nombre de vulnérabilités exploitées non corrigées après sept
          jours.
    \item La revue mensuelle des règles de pare-feu.
\end{enumerate}
\end{boxexo}

\begin{boxexo}[title={\ding{45}\quad Exercice~13.2. Compléter la matrice \niveaumoyen}]
\begin{enumerate}[(a)]
    \item Ajoutez à la matrice (Tableau~13.3) le contrôle «~protection
          contre le déni de service chez le fournisseur d'accès~» et le
          risque «~règle de filtrage temporaire maintenue~». Reliez-les.
    \item Pour chaque risque, dites s'il a au moins un contrôle \OUI{}.
          Lesquels ne sont pas maîtrisés~?
    \item Calculez $\mu$, $U$ et $H$ des neuf contrôles du tableau, puis
          celles des seuls contrôles qui traitent R-AG-01. Que montre la
          comparaison~?
\end{enumerate}
\end{boxexo}

\begin{boxexo}[title={\ding{45}\quad Exercice~13.3. Des indicateurs pour les prestataires \niveaumoyen}]
\begin{enumerate}[(a)]
    \item Proposez deux indicateurs de risque pour R-DC-02 (DIGICORE), avec
          leurs seuils et leur source.
    \item Proposez un indicateur de risque pour R-AG-02 (GAB) une fois le
          remplacement achevé.
    \item Dites, pour chacun, ce que déclenche le passage au rouge, et qui
          décide.
\end{enumerate}
\end{boxexo}

\begin{boxexo}[title={\ding{45}\quad Exercice~13.4. La page de l'ANPP \niveauexpert}]
À partir des Chapitres~\ref{chap:referentiels} et~\ref{chap:droit},
préparez le tableau de bord GRC d'une page destiné au conseil
d'administration de l'ANPP.
\begin{enumerate}[(a)]
    \item Construisez le registre des obligations (au moins six lignes),
          avec leur statut.
    \item Calculez $\mu$, $U$ et $H$ de ce registre et situez-le dans le
          quadrant.
    \item Rédigez les six rubriques de la page, sur le modèle de celle de
          M.~ATANGANA, dont la rubrique «~ce que nous ne savons pas~».
    \item Proposez deux tests automatiques quotidiens réalistes pour
          l'ANPP, et dites comment l'audit les validera.
\end{enumerate}
\end{boxexo}

\bigskip

\begin{boxinfo}
\textbf{Pour aller plus loin.}
\begin{itemize}
    \item L'ISO~31000~\cite{iso31000_2018}, pour l'intégration de la gestion
          des risques dans la gouvernance.
    \item Le cadre COSO de gestion des risques~\cite{coso2017erm}, pour
          l'appétence et les indicateurs.
    \item Le NIST CSF~2.0~\cite{nist2024csf}, fonction \emph{Gouverner},
          catégorie de supervision.
\end{itemize}
\end{boxinfo}

\begin{boxinfo}
\textbf{Transition.} Les chapitres de la Partie~IV ont traité la pratique
par morceaux~: le SMSI, l'audit technique, la GRC. Le
Chapitre~\ref{chap:cas-integrateurs} les assemble dans trois cas complets,
de l'incident au rapport.
\end{boxinfo}
